\documentclass[11pt]{article}

\usepackage[margin=1in]{geometry}
\usepackage{amsmath,amssymb,amsthm}
\usepackage{booktabs,array}
\usepackage{enumitem}
\usepackage{xcolor}
\usepackage[colorlinks=true,linkcolor=blue!50!black,citecolor=blue!50!black]{hyperref}

% Drafting helpers: delete calls as you fill in
\newcommand{\todo}[1]{\par\smallskip\noindent\textcolor{red!70!black}{\textbf{[FILL:}~#1\textbf{]}}\par\smallskip}
\newcommand{\decide}[1]{\par\smallskip\noindent\textcolor{orange!80!black}{\textbf{[DECIDE:}~#1\textbf{]}}\par\smallskip}

% Notation
\newcommand{\Gd}{G_{d}}          % those in power
\newcommand{\Gn}{G_{n}}          % general public
\newcommand{\Gr}{G_{r}}          % resistant subgroup (the movement)
\newcommand{\sd}{\sigma_{s}}     % dominant (silence) schema
\newcommand{\sr}{\sigma_{r}}     % resistant (disclosure) schema
\newcommand{\E}{\mathbb{E}}

\theoremstyle{definition}
\newtheorem{definition}{Definition}[section]
\theoremstyle{remark}
\newtheorem{remark}{Remark}[section]

\title{Feedback, Reputation and Status in a Resistant Social Reward Economy}
\author{Aaron Avram}
\date{September 29, 2026}

\begin{document}
\maketitle

% ============================================================
\section{Introduction}
% ============================================================

\subsection{Motivation and gap}
\todo{One paragraph. The SRE paper describes resistance as an alternative reward economy but leaves feedback, reputation and status loosely specified, and uses feedback quantities in its inequalities without defining who gives feedback to whom or how it is weighted. This report supplies those three definitions.}
The \textit{Social Reward Economies} paper describes resistance as an alternative reward economy
but leaves feedback, reputation and status loosely specified with feedback
quantities used in the inequalities measuring the effectiveness of resistance. This report
provides these three definitions.

\subsection{Scope}
\todo{State plainly what this report does and does not do. It defines three quantities and shows how they express the features of resistance in an example. It does not define a decision rule, an attention parameter, success conditions for resistance, gossip structure or any dynamics. Give one sentence on why: those depend on the definitions here and should follow their review.}
This report defines three quantities and how they express features of resistance in the \#MeToo example.
It does not define a decision rule, success conditions for resistance and other key parts
of implementing resistance in a multi-agent system. The reason for this is that these depend on the definitions
of the three stated quantities and should not be considered until these definitions are finalized.

% ============================================================
\section{The Example}
% ============================================================

\subsection{\#MeToo}
\todo{Three types of people: those in power ($\Gd$), the general public ($\Gn$), the movement ($\Gr$). Keep the observation from the earlier report that the public is not strongly affected by either side yet is the majority. Describe the pre-movement norm (silence) and the movement's counter-norm (disclose and hold accountable).}
An example is the formation of the \#MeToo Movement. There are three types of people in this scenario:
those in charge in the dominant group (such as high ranking corporate officials)
the members of the \#MeToo movement and finally the general public (which is just everyone else in society). In this setup
it is clear that those in charge benefit at the expense of many women in the workplace
and by extension many members of the \#MeToo movement. Hence, the group in charge
will have an incentive to preserve the dominant norm, which advocates for not calling out sexual harassment,
as this directly benefits them. Conversely, women and the \#MeToo movement wish to overturn the dominant
norm in favour of their resistant one, that sexual harassment should be made public and the perpetrators
should be shamed, as this benefits them. However, the general public is not strongly affected by the behaviours
of either of the above groups, yet they make up a majority of the population.

The pre-movement norm is that individuals should stay silent about sexual assault and the movement's counter-norm
is that they should instead disclose these cases and hold the perpetrators accountable.

\subsection{Schemas and states}
\decide{Use two schemas, dominant $\sd$ (silence / excuse) and resistant $\sr$ (disclose / hold accountable), folding the earlier report's reformed and resistant schemas together? Nothing in this report needs three.}
\todo{Let $S_r \subseteq S$ be the states where sexual harassment is at issue. Outside $S_r$, utilities are roughly uniform.}

There are two schemas: dominant $\sd$ (silence/excuse) and resistant $\sigma_r$ (disclose/hold accountable)
Let $S_r \subseteq S$ be the set of states where sexual harassment is at issue. Outside
of $S_r$, utilities are roughly uniform.

\subsection{Utilities}
\todo{Reuse the earlier report's utility distribution. It also serves as the feedback table, since feedback is derived from utilities (Section~\ref{sec:feedback}).}
\begin{table}[h]
\centering
\begin{tabular}{lcc}
\toprule
Type of agent $j$ & $u_j(s,\sd)$ & $u_j(s,\sr)$ \\
\midrule
$\Gd$ (in power)  & $\gg 0$ & $< 0$ \\
$\Gn$ (public)    & $> 0$   & $> 0$ \\
$\Gr$ (movement)  & $\ll 0$ & $\gg 0$ \\
\bottomrule
\end{tabular}
\caption{Utilities for $s \in S_r$. \todo{Confirm entries against the earlier report.}}
\end{table}

\subsection{Contrast cases}
\todo{One or two sentences each. Accessibility: the public and those in charge do not gain from the old norm. Liverpool fans: a subgroup with its own norm and network, unaffected by and not affecting the city. Used only in Section~\ref{sec:contrast}.}
The above example is only a specific case where the dominant group and the subgroup are in direct opposition.
However, there are other cases. First, consider the example of the accessibility
movement: the public and those in charge do not gain from the old norm. Or,
consider fans of a sports club: this is a subgroup with its own norm and network,
unaffected by and not affecting the rest of society.

% ============================================================
\section{Minimal Setting}
% ============================================================
\todo{Notation only, as short as possible. Do not introduce anything that is not used in Section~\ref{sec:defs}.}
Some notation:

\begin{itemize}[nosep]
  \item States $s \in S$ with probabilities $p(s)$; actions $A(s)$.
  \item Utilities $u_j(s,a) \in \mathbb{R}$, heterogeneous across agents, positive or negative.
  \item Agent $k$'s behaviour as a distribution over actions, $O_k(s)$: the softmax policy from Learning Norms. This differs from the SRE paper's $O_i(s)$, which denotes a schema
  and can be viewed as a proxy for $k$'s actual schema.
  \item Groups: a larger group $G$ and a subgroup $\Gr \subset G$. A generic group is written $H \in \{G, \Gr\}$.
  \item Participation: rate $\mu_{j,p}^H$ and activity probability $\theta(\mu) = 1 - e^{-\mu}$, used only to define audience weights.
\end{itemize}

% ============================================================
\section{Definitions}\label{sec:defs}
% ============================================================
\todo{Use the same three parts for each definition: (i) the definition; (ii) the simpler alternative and what it cannot represent; (iii) the \#MeToo instance.}
\subsection{Social feedback}\label{sec:feedback}
\begin{definition}[Social feedback]
For agents $j \neq k$ and a state $s$, the feedback from $j$ to $k$ is $j$'s expected benefit from $k$'s behaviour:
\begin{equation}
  f_{jk}(s) \;=\; \E_{a \sim O_k(s)}\big[\, u_j(s,a) \,\big].
\end{equation}
\end{definition}
\todo{Simpler alternative: an unsigned endorsement (as in the Learning Norms status). It cannot express sanction. \#MeToo instance: a person in power gives negative feedback to a survivor who discloses; movement members give positive feedback. Note that feedback is defined between any two agents, in or out of the same group.}
Here it is crucial that $u_j(s, a)$ can be both positive and negative so that social feedback
can be both positive and negative. This is necessary as in the \#MeToo example:
a person in power gives negative feedback to a surivor who discloses but the movement members
give this action positive feedback.

\begin{remark}
\todo{Point out the indexing: $f_{jk}$ uses $j$'s utility $u_j$ evaluated on $k$'s behaviour. Using $u_k$ instead would give $k$'s own utility, which is personal utility, not feedback.}
\end{remark}

\subsection{Audience weights (auxiliary)}
\begin{definition}[Audience weights]
For a group $H$, let
\begin{equation}
  \nu_j^H \;=\; \frac{\theta(\mu_{j,p}^H)}{\sum_{l \in H} \theta(\mu_{l,p}^H)}, \qquad j \in H.
\end{equation}
\end{definition}
\todo{Explain in two sentences why $\theta$ is normalised (it is an activity probability, not a distribution) and why participant, not actor, rates are used. This is an auxiliary object, not a main definition.}
This normalization is important so that we get a probability distribution over $H$ which expresses
the probability of any given agent participating. This is important as it relates
to how much other agents care about $j$'s feedback. If $j$ does not participate regularly
then their feedback is less valuable.

\subsection{Reputation}
\begin{definition}[Reputation]
For a group $H$ and a member $k \in H$,
\begin{equation}
  R_k^H(s) \;=\; \sum_{j \in H \setminus \{k\}} \nu_j^H \, f_{jk}(s).
\end{equation}
\end{definition}
\todo{Reputation is the audience-averaged feedback within $H$. Simpler alternative: one global reputation, which averages the two audiences together and cannot represent a counterpublic. \#MeToo instance: an advocate has high $R^{\Gr}$ and a much lower, possibly negative, $R^{G}$; a protected person in power has high $R^{G}$.}
\decide{Should reputation be defined for non-members? Your constraint is that members of $G$ outside $\Gr$ have no reputation in $\Gr$, so the definition above is restricted to $k \in H$. The cost: a person in power has no $R^{\Gr}$, and the movement's rejection of them shows up as \emph{negative feedback from $\Gr$ members} (Section~\ref{sec:feedback}), not as a negative reputation. This differs from the earlier report, which spoke of the abuser's ``highly negative reputation'' in the movement. Alternatively, define $R^{\Gr}_k$ for all $k \in G$ using raters in $\Gr$, and say that only the \emph{rater} role is closed to outsiders.}

Reputation is the audience-averaged feedback within $H$. It is important that there are separate reputation structures
in $G$ and $G_r$ as otherwise there could not be any counter-public. In the \#MeToo example,
an advocate of the movement has high $R^{G_r}$ and a much lower, possibly negative $R^G$
while a person in power has high $R^G$.

\subsection{Status}
\begin{definition}[Status]
For $k \in H$, let $w_j^H = \nu_j^H\,[R_j^H]_+$ be rater $j$'s weight, where $[x]_+ = \max(x,0)$, and set
\begin{equation}
  S_k^H(s) \;=\; \sum_{j \in H \setminus \{k\}} w_j^H \, f_{jk}(s).
\end{equation}
\end{definition}
\todo{Status is reputation-weighted feedback: feedback from well-regarded raters counts for more. Why clip at zero: a discredited rater is discounted rather than having his feedback sign-flipped. \#MeToo instance: an early discloser receives negative feedback from the heavily weighted powerful in $G$ (low or negative $S^{G}$) but strongly positive feedback from respected movement members ($S^{\Gr}$ high). This is the support that makes speaking out survivable.}
\decide{Clipping at zero versus another monotone nonnegative map for the weights.}

Status is reputation-weighted feedback: feedback from well-regarded raters counts more.
The reputation term in the feedback weight is cliped at zero as a discredited rater
should be ignored not have their feedback sign-flipped. In the \#MeToo example,
an early discloser recieves negative feedback from powerful members of $G$
but strongly positive feedback from respected movement members. This is the support
that makes speaking out survivable.

\begin{remark}
This is not the status defined in Learning Norms, where $S_k$ sums each follower's reputation $R_i(\pi_k)$, i.e.\ the benefit to everyone else. This status sums raters' own perceived benefit; it recovers that definition only if the weights are follower indicators and $f$ is replaced by $R_i$. 
This setup also has a hierarchy of influence among the defined quantities: weights depend on reputation and status depends on weights.
\end{remark}

% ============================================================
\section{How the Three Quantities Pertain to Resistance}
% ============================================================
\todo{Frame this section as description, not as a result: which pattern across feedback, reputation and status distinguishes a resistant subgroup from an ordinary one. Do not state success conditions, thresholds or decision rules.}

\subsection{The configuration before the movement}
\todo{Fill the table below from the utility table and the definitions. Suggested reading: the powerful and the public give positive feedback to behaviour $\sd$ and negative or muted feedback to $\sr$; reputation and status in $G$ are high for those in power and for those who stay silent.}
\begin{table}[h]
\centering
\begin{tabular}{lccc}
\toprule
Agent (behaviour) & $R^{G}$ & $S^{G}$ & Feedback from $\Gr$ \\
\midrule
Person in power ($\sd$)    & & & \\
Advocate who discloses ($\sr$) & & & \\
Bystander ($\sd$)          & & & \\
\bottomrule
\end{tabular}
\caption{Standing in $G$ for representative agents. \todo{Verify each sign against the definitions and the utility table; state assumptions about weights.}}
\end{table}

\subsection{An alternative economy: the pattern that signals resistance}
\todo{Describe, in words and one table, what it means for $\Gr$ to have its own economy in terms of the three quantities. Suggested content: for the same behaviour, the sign of feedback given by members of $\Gr$ differs from the sign given by the rest of $G$; the ranking of agents by reputation within $\Gr$ differs from the ranking in $G$; and status within $\Gr$ is held by different agents than status in $G$.}
\begin{table}[h]
\centering
\begin{tabular}{lcc}
\toprule
Agent (behaviour) & $R^{\Gr}$, $S^{\Gr}$ (movement's view) & $R^{G}$, $S^{G}$ (larger group's view) \\
\midrule
Advocate who discloses ($\sr$) & & \\
Movement leader ($\sr$)        & & \\
\bottomrule
\end{tabular}
\caption{The same agents evaluated by the two audiences. \todo{The point of the table is the sign reversal between columns.}}
\end{table}
\begin{remark}
\todo{Say that this is a description of what an alternative economy looks like, and that the SRE paper's Eqs.\ (1)--(2) are stated in terms of weighted feedback of exactly this kind. Say that how such feedback determines whether the subgroup's norm spreads is left for later.}
\end{remark}

\subsection{The cost of speaking out}
\todo{Qualitative only. A discloser receives negative feedback and low status from the larger group but positive feedback and high status inside $\Gr$. Compare $S^{\Gr}$ and $S^{G}$ for a discloser to say why speaking out is costly and why the subgroup's status is what sustains it. Do not define a threshold.}

\subsection{Asymmetry between the groups}
\todo{Members of $G$ outside $\Gr$ give no weighted feedback inside $\Gr$ (the constraint from Section 4), while members of $\Gr$ remain part of $G$ and are evaluated there. Say why this matters for the example: the subgroup can be affected by outsiders' feedback, but outsiders are not affected by the subgroup's internal ranking.}
\decide{Do members of $\Gr$ carry reputation and status in $G$? Here they do, which is what makes ``high in $\Gr$, low in $G$'' expressible.}

\subsection{Contrast cases}\label{sec:contrast}
\todo{Show that the same three quantities describe the other examples differently.}
\begin{table}[h]
\centering
\begin{tabular}{p{3cm}p{5cm}p{6cm}}
\toprule
Example & Feedback from the wider group to the subgroup's norm & Reputation and status pattern \\
\midrule
\#MeToo         & & \\
Accessibility   & & \\
Liverpool fans  & & \\
\bottomrule
\end{tabular}
\caption{Pattern of the three quantities in each example.}
\end{table}

% ============================================================
\section{Discussion}
% ============================================================

\subsection{Relation to Learning Common Norms}
\todo{Restate the status difference. One sentence on what carries over (signed vs.\ unsigned is new; utilities and softmax policies carry over).}

\subsection{What is deliberately left out}
\todo{One short list, no detail: how feedback aggregates into a choice between schemas; how much attention agents give each group; formal conditions for a subgroup norm to spread or to displace the dominant one; who interacts with whom; how any of these quantities change over time. State that these are the natural next step once the three definitions are settled.}

\subsection{Limitations}
\begin{itemize}[nosep]
  \item \todo{The example is used to motivate the definitions, so agreement with the example is plausibility, not validation.}
  \item \todo{The definitions are static and say nothing about whether the described patterns arise.}
  \item \todo{Dependence on the choices marked above (schema count, weight map, reputation for non-members).}
\end{itemize}

\end{document}